\section{Introduction}\label{sec:intro}

The phase retrieval problem asks whether a vector or function can be recovered
from phaseless measurements. In the real setting this means recovery up to
global sign. The finite-dimensional theory is by now extensive; see
\cite{bandeira2014saving, eldar2014stability} and the survey
\cite{grohs2020survey}. In infinite dimensions the situation is more
delicate: phase retrieval can hold in infinite-dimensional Hilbert spaces
\cite{cahill2016infinite}, but uniform stability may fail for continuous
frames \cite{alaifari2017continuous}, and natural measurement models such as
Gabor frames may be severely ill-posed \cite{alaifari2021gabor}.

We adopt the function-space point of view from \cite{FOTP}. Given a closed
subspace $E$ of a real function space $X$, we say that $E$ does \emph{stable
phase retrieval} with constant $C$ if
\[
\min_{\sigma \in \{\pm 1\}} \norm{f - \sigma g}_X
\le C \, \bigl\| \abs{f} - \abs{g} \bigr\|_X
\qquad
\text{for all } f, g \in E.
\]
The map $f \mapsto \abs{f}$ loses sign information, so the left-hand side
optimizes over a global sign flip; the inequality says that this is the only
obstruction to recovery. The main results of \cite{FOTP} show that many
order-continuous Banach lattices admit no infinite-dimensional subspace with
this property, while \cite{christ2022examples, FOTP} provide positive
examples in $L^2$-spaces. The orthogonal reduction of \cite{FOTP} is the
starting point of our argument: to verify stable phase retrieval on a
subspace, it suffices to check it for orthogonal pairs in the same
two-dimensional span.

In \cite{calderbank2022stable}, Calderbank, Daubechies, Freeman, and Freeman
proved stable phase retrieval for a model involving shifts and proposed the
following conjecture: if $(\xi_i)_{i \ge 1}$ is an orthonormal sequence of
independent mean-zero random variables in $L^2$, then
$\overline{\Span}(\xi_i)$ does stable phase retrieval if and only if there
exist $0 < A \le B < 1$ with
\[
A \le \norm{\xi_i}_{L^1} \le B
\]
for all but at most one $i$. The lower bound prevents $\xi_i$ from
concentrating near zero; the upper bound prevents it from being
Rademacher-valued (in which case $\abs{\xi_i}$ is constant and carries no
phase information). Together they ensure that each coordinate contributes
usable information about the signs of the coefficients.

Our main result proves the sufficiency direction of this conjecture with
a fully explicit stability constant.

\subsection{The main theorem}

Let $\xi_1, \dots, \xi_n$ be independent real-valued random variables on a
probability space $(\Omega, \mu)$, each satisfying
\[
\E[\xi_i] = 0,
\qquad
\E[\xi_i^2] = 1,
\qquad
A \le \E[\abs{\xi_i}] \le B
\]
for fixed $0 < A \le B < 1$. Given unit vectors $a, b \in \R^n$ with
$\ip{a}{b} = 0$, form the random linear combinations
$X_a = \sum_i a_i \xi_i$ and $X_b = \sum_i b_i \xi_i$. The \emph{modulus gap}
\[
\delta(a,b) = \sqrt{\E\bigl[(\abs{X_a} - \abs{X_b})^2\bigr]}
\]
measures how distinguishable $X_a$ and $X_b$ are from their absolute values
alone. After the orthogonal reduction, stable phase retrieval with constant $C$
is equivalent to $\delta(a,b) \ge \sqrt{2}/C$ for all orthogonal unit pairs.

\begin{theorem}[Quantitative orthogonal lower bound]\label{thm:main}
\leanname{Showcase\.quantitative\_orthogonal\_lower\_bound}
There exists a constant $c_\perp(A,B) > 0$, which can be chosen to be
\[
c_\perp(A,B)
= \tfrac{1}{2}
  \min\bigl(c_{\mathrm{cross}}(A),\;
  c_{\mathrm{line}}(A,B)\bigr),
\]
such that $c_\perp(A,B) \le \delta(a,b)$
for every pair of orthogonal unit vectors $a, b \in \R^n$.
The constant depends only on $A$ and $B$, not on the dimension $n$ or the
laws of $\xi_i$. The quantities $c_{\mathrm{cross}}$ and
$c_{\mathrm{line}}$ are defined in Section~\ref{sec:constants}.
\end{theorem}

\begin{remark}
The companion paper \cite{compactness} proves the same conclusion (without
explicit constants) by a compactness argument that also allows one
exceptional variable with no $L^1$ bound. The present proof makes each step
of that argument quantitative.
\end{remark}

\subsection{Proof strategy}\label{sec:strategy}

Rather than lower-bounding $\delta$ directly, we work with the auxiliary
quantity
\[
\eps(a,b) = \E\bigl[\abs{X_a^2 - X_b^2}\bigr],
\]
which is more amenable to algebraic manipulation. A Cauchy--Schwarz argument
(Theorem~\ref{thm:eps-le-modulus-app}) gives $\eps \le 2\delta$, so any lower
bound on $\eps$ transfers to $\delta$ at the cost of a factor of two.

The proof splits the coordinates into a \emph{head} (the few coordinates
where $a$ or $b$ has a large coefficient) and a \emph{tail} (the remaining
coordinates, each individually small). A sign dichotomy on the tail
(Lemma~\ref{lem:tail-dichotomy}) then determines which of two branches
applies:

\begin{enumerate}[label=(\roman*)]
\item \textbf{Concentrated branch} (Section~\ref{sec:head}).
For some sign $\tau = \pm 1$, the tail of $a - \tau b$ is small, so $a$ and
$\tau b$ nearly agree outside the head. Bounded probe functions extract
entries of the head difference matrix from the observable
$X_a^2 - X_b^2$, each entry being controlled by $K(A,B) \cdot \eps$.
The orthogonality $\ip{a}{b} = 0$ forces this matrix to have large Frobenius
norm, giving $\eps \ge c_{\mathrm{line}}(A,B)$.

\item \textbf{Diffuse branch} (Section~\ref{sec:blocks}).
Both $a + b$ and $a - b$ have substantial tail energy, spread across many
small coefficients. A greedy block partition, symmetrization via the product
space, Sperner's antichain bound, and a Chebyshev counting argument
together show that $X_a$ and $X_b$ are both unlikely to be near zero,
giving $\eps \ge c_{\mathrm{cross}}(A)$.
\end{enumerate}

\noindent
In both cases $\eps \ge \min(c_{\mathrm{cross}}, c_{\mathrm{line}})$, and
Theorem~\ref{thm:main} follows.


\section{Setup: the head--tail decomposition}\label{sec:setup}

Throughout, $a, b \in \R^n$ are unit vectors with $\ip{a}{b} = 0$, and
$\xi_1, \dots, \xi_n$ are independent random variables satisfying the moment
conditions of Section~\ref{sec:intro}. We write $X_c = \sum_i c_i \xi_i$ for
any coefficient vector $c \in \R^n$.

The first step is to separate the coordinates where $a$ and $b$ have large
coefficients from those where they are small. This is the quantitative
counterpart of the profile decomposition in \cite{compactness}: the large
coordinates form a finite set on which we can extract detailed information
via probes, while the small coordinates contribute a ``diffuse'' component
that is controlled by anticoncentration.

\begin{definition}[Head set, restriction, and tail]\label{def:head-set}
\leanname{Core\.FiniteModel\.headSetFin}
The \emph{head set} is
$H = H(A,a,b)
= \bigl\{i : \max(\abs{a_i}, \abs{b_i}) > \lambda(A)\bigr\}$.
For $S \subseteq \{1, \dots, n\}$ and $c \in \R^n$, write $\restr_S c$ for
the vector that agrees with $c$ on $S$ and vanishes elsewhere
(\leanref{restrictCoeff}), and
$\tail_S c = c - \restr_S c$ for its complement
(\leanref{tailCoeff}).
\end{definition}

Since each $i \in H$ contributes at least $\lambda(A)^2$ to
$a_i^2 + b_i^2 \le 2$, the head set is finite with
$\card{H} \le 2/\lambda(A)^2 \le d^*(A)$
(\leanref{headSetFin\_card\_le\_dStar}).
Conversely, every tail coordinate satisfies
$\max(\abs{(\tail_H a)_i}, \abs{(\tail_H b)_i}) \le \lambda(A)$
by construction
(\leanref{headSetFin\_tail\_max\_le\_lambda}).

\label{thm:head-card}%
\label{thm:tail-bound}%

\subsection{The head difference matrix}

The head difference matrix captures the discrepancy between $a$ and $b$ on
the large coordinates. It is a rank-two symmetric matrix whose Frobenius norm
is forced to be large by the orthogonality constraint whenever the tails
are small.

\begin{definition}[Difference matrix]\label{def:diff-matrix}
\leanname{Core\.FiniteModel\.finDH}
The head difference matrix $D_H(a,b)$ is the symmetric matrix with entries
\[
D_H(a,b)_{ij}
= (\restr_H a)_i (\restr_H a)_j
- (\restr_H b)_i (\restr_H b)_j
\]
and Frobenius norm
$\norm{D_H}_F = \sqrt{\sum_{i,j} D_H(a,b)_{ij}^2}$.
\end{definition}

\subsection{The epsilon--modulus bridge}

The two key quantities measuring phase-retrieval distinguishability are
\[
\eps(a,b) = \E\bigl[\abs{X_a^2 - X_b^2}\bigr],
\qquad
\delta(a,b)
= \sqrt{\E\bigl[(\abs{X_a} - \abs{X_b})^2\bigr]}.
\]
Since
$\abs{X_a^2 - X_b^2}
= \abs{X_a + X_b} \cdot \abs{X_a - X_b}$,
Cauchy--Schwarz and the variance identity
$\E[X_c^2] = \sum c_i^2$
(\leanref{Imported\.L2Theory\.randL2Sq\_finLinComb\_eq})
yield $\eps \le 2\delta$
(Theorem~\ref{thm:eps-le-modulus-app}). The rest of the
proof is devoted to showing $\eps$ is bounded below.


\subsection{Table of constants}\label{sec:constants}

All constants depend only on $A$ and $B$ and are explicitly computable.
The probe bound, gluing threshold, and head/tail cutoff are
\[
K = \max\!\bigl(\tfrac{4}{A^2},\,\tfrac{2}{1-B}\bigr),
\qquad
\theta = \tfrac{A^2}{8},
\qquad
\lambda = \tfrac{A^3}{2^{20}}.
\]
The head set has at most $d^* = 2 + 2^{41}/A^6$ elements, and the
diffuse threshold is $s_0 = A^4/2^{20}$. The two branch constants and
the stability constant are
\[
c_{\mathrm{cross}} = \tfrac{A^8}{2^{57}},
\qquad
c_{\mathrm{line}} = \tfrac{A^2}{2^8\, d^*\! K},
\qquad
c_\perp = \tfrac{1}{2}
\min(c_{\mathrm{cross}}, c_{\mathrm{line}}).
\]
The block partition in Section~\ref{sec:blocks} also uses
\[
M = \lceil 2^{18}/A^2 \rceil,
\quad
\eta(s) = \tfrac{s}{4M},
\quad
\lambda_0(s) = \tfrac{A\sqrt{s}}{2048},
\quad
r_0(s) = \tfrac{A^2\sqrt{s}}{65536}.
\]\label{sec:constants-end}
